\documentclass[9pt,t,serif,professionalfont,xcolor={table,dvipsnames},noamsthm]{beamer}

\usepackage[perpage,symbol*,para,multiple]{footmisc}
% Getting undefined commands in aux file
\makeatletter
\newcommand\FN@pp@footnote@aux[2]{\relax}
\makeatother
\PassOptionsToPackage{pdfpagemode=FullScreen}{hyperref}
\setbeamersize{text margin left = 0.5cm,
               text margin right = 0.3cm}
\mode<presentation>
{
  \usetheme{boxes}
  \setbeamercovered{invisible}
  \setbeamertemplate{navigation symbols}{} 
}
\definecolor{alertcolor}{RGB}{203, 0 , 34}
\setbeamercolor{alerted text}{fg=alertcolor}

% set the fonts
\usefonttheme{professionalfonts} % using non standard fonts for beamer
\usefonttheme{serif} % default family is serif
\usepackage{../presentation,../presentationfonts,../dash}

% load the pdfanim style, should be done after hyperref
% \usepackage{pdfanim}

% load and initialise animation
% arguments:
%   \PDFAnimLoad[options]{name}{xxx}{number}
%   - options: in this example the animation is looped
%   - name: name of the animation
%     (with this name it can be reused several times in the document)
%   - number: number of animation frames / files (n)
% the animation will be composed of files
% xxx0.pdf, xxx1.pdf ... xxx(n-1).pdf


\title{Preparatory material}

\author{Jim Hef{}feron}
\institute{
  St~Michael's College \\
  and University of Vermont\\[1ex]
  \texttt{hefferon.net}
}
\date{}

% This is only inserted into the PDF information catalog.
\subject{Introduction to Computability}

% \usepackage{catchfilebetweentags} % read text from background.tex
\usepackage{filebetweentags}
% \newcommand{\catchfilefn}{../../prologue/prologue.tex}

% 2026-Jun-30 JH This package has stopped working.  I fadged up a hack
% using a Python file. 
% % Grab the text from another file, between tags %<*tagstring> and %</tagstring>
% % #1 file name, including the ".tex" 
% % #2 tag string
% \newcommand{\ExecuteMetaData}[2][\catchfilefn]{%
%   \input{|python3 catchfilebetweentags.py -D -f #1 -t #2 -s} % -D fordebugging 
%   }

% % Keep from swallowing end-of-lines
% %   https://tex.stackexchange.com/a/40704/121234
% \usepackage{etoolbox}
% \makeatletter
% \patchcmd{\CatchFBT@Fin@l}{\endlinechar\m@ne}{}
%   {}{\typeout{Unsuccessful patch!}}
% \makeatother

% Get cross refs from .aux files.  Must load after zref-clever, which
% is loaded in presentation.sty
\usepackage{zref-xr}
  %\zexternaldocument{../../book}
  \zexternaldocument{../../appendix/strings}[../../book.pdf]
  % \zexternaldocument{../../appendix/functions}[../../book.pdf]
  % \zexternaldocument{../../appendix/proplogic}[../../book.pdf]

% Directory of this chapter's asy material
\newcommand{\asydir}{../../appendix/asy/}
\renewcommand{\tmexercises}{\asydir tm_exercises/pdfs/}
% Directory of the src/asy material
\newcommand{\generalasydir}{../../asy/}

\usepackage{multirow,bigstrut}

\usepackage{pythontex}
\begin{document}
\begin{frame}
  \titlepage
\end{frame}



\section{Strings}
% ------------------------------------------------------
\renewcommand{\catchfilefn}{../../appendix/strings.tex}
\begin{frame}
  \frametitle{Alphabets and strings}
\FileBetweenTags[\catchfilefn]{def:Alphabet}
\FileBetweenTags[\catchfilefn]{def:String}

\FileBetweenTags[\catchfilefn]{discussion:SymbolsMayContainMoreThanOneChar}

\pause
\FileBetweenTags[\catchfilefn]{discussion:AlphSymbol}

\FileBetweenTags[\catchfilefn]{def:StringLength}

\pause
\FileBetweenTags[\catchfilefn]{dis:OmitCommas}
\end{frame}


\begin{frame}
\FileBetweenTags[\catchfilefn]{def:Bitstrings}

\FileBetweenTags[\catchfilefn]{def:KleeneStar}

\pause
\FileBetweenTags[\catchfilefn]{def:StringOps}

\FileBetweenTags[\catchfilefn]{def:Palindrome}
\end{frame}




\section{Functions}
% ------------------------------------------------------
\renewcommand{\catchfilefn}{../../appendix/functions.tex}

\begin{frame}
  \frametitle{Definition}

\FileBetweenTags[\catchfilefn]{def:functioni}
\FileBetweenTags[\catchfilefn]{def:functionii}

\pause
\FileBetweenTags[\catchfilefn]{def:functioniii}
\end{frame}

\begin{frame}
  \frametitle{Showing a function}
  
\FileBetweenTags[\catchfilefn]{disc:Picturesi}
\end{frame}

\begin{frame}
\FileBetweenTags[\catchfilefn]{disc:Picturesii}
\end{frame}


\begin{frame}
  \frametitle{Range}
\FileBetweenTags[\catchfilefn]{def:Range}
The range of the logarithm function $\map{\ln}{\R^+}{\R}$ is the entire
set of reals. 
\end{frame}

\begin{frame}
  \frametitle{Partial functions}
\FileBetweenTags[\catchfilefn]{def:PartialFunctions}
\end{frame}


\begin{frame}
  \frametitle{Showing that something is well-defined}
\FileBetweenTags[\catchfilefn]{def:WellDefinedi}

\FileBetweenTags[\catchfilefn]{def:WellDefinedii}

\vspace{10pt}
\begin{center}
  \includegraphics{asy/functions000.pdf}
\end{center}
\end{frame}

\begin{frame}
\FileBetweenTags[\catchfilefn]{def:WellDefinediii}  
\end{frame}

\begin{frame}
\FileBetweenTags[\catchfilefn]{def:WellDefinediv}  
\end{frame}

\begin{frame}
\FileBetweenTags[\catchfilefn]{def:WellDefinedv}  
\end{frame}



\begin{frame}
  \frametitle{One-to-one and onto}
\FileBetweenTags[\catchfilefn]{def:OneToOne}  
\end{frame}

\begin{frame}
\FileBetweenTags[\catchfilefn]{def:Onto}  
\end{frame}


\begin{frame}
  \frametitle{Correspondence}
\FileBetweenTags[\catchfilefn]{def:Correspondence}  
\end{frame}


\begin{frame}
  \frametitle{Function composition}
\FileBetweenTags[\catchfilefn]{def:Compositioni}  

\pause
\FileBetweenTags[\catchfilefn]{def:Compositionii}  
\end{frame}


\begin{frame}
\FileBetweenTags[\catchfilefn]{def:IdentityFunction}
    
\pause
\FileBetweenTags[\catchfilefn]{def:Inverse}
\end{frame}



\section{Propositional logic}
% ------------------------------------------------------
\renewcommand{\catchfilefn}{../../appendix/proplogic.tex}

\begin{frame}
  \frametitle{Propositions}
\FileBetweenTags[\catchfilefn]{def:Proposition}

\pause
\FileBetweenTags[\catchfilefn]{def:Operations}

\FileBetweenTags[\catchfilefn]{def:BooleanVariables}
\end{frame}

\begin{frame}
\FileBetweenTags[\catchfilefn]{def:TruthTables}
\end{frame}


\begin{frame}
  \frametitle{Logical operators}
\FileBetweenTags[\catchfilefn]{def:Operators}

\pause
\FileBetweenTags[\catchfilefn]{def:Equivalence}
\end{frame}


\begin{frame}
  \frametitle{Boolean Algebra}
\FileBetweenTags[\catchfilefn]{def:BooleanAlgebra}
\end{frame}


\begin{frame}
  \frametitle{DNF and CNF}
Given a truth table, we can produce an expression with
that table.
We will give two such expression types.
For the first we use these tables as examples.
\begin{center}
\FileBetweenTags[\catchfilefn]{def:DNFi}
\hspace*{2em}
\FileBetweenTags[\catchfilefn]{def:DNFii}
\end{center}
\FileBetweenTags[\catchfilefn]{def:DNFiii}
\FileBetweenTags[\catchfilefn]{def:DNFiv}
\end{frame}

\begin{frame}
Here is that ($*$) expression again.
\FileBetweenTags[\catchfilefn]{def:DNFiv}%
\FileBetweenTags[\catchfilefn]{def:DNFv}
\end{frame}

\begin{frame}
\FileBetweenTags[\catchfilefn]{def:CNFi}
Here is the example table again.
\begin{center}
  \FileBetweenTags[\catchfilefn]{def:DNFii}
\end{center}
\FileBetweenTags[\catchfilefn]{def:CNFii} 

\FileBetweenTags[\catchfilefn]{def:CNFiii} 
\end{frame}




\end{document}

